\documentclass[11pt]{article}

\usepackage{amsmath,amsthm,amssymb}
\usepackage[margin=1in]{geometry}

\newtheorem{theorem}{Theorem}
\newtheorem{lemma}{Lemma}
\newtheorem{proposition}{Proposition}
\newtheorem{corollary}{Corollary}
\theoremstyle{definition}
\newtheorem{remark}{Remark}
\newtheorem{definition}{Definition}

\DeclareMathOperator{\Des}{Des}
\DeclareMathOperator{\maj}{maj}
\DeclareMathOperator{\comaj}{comaj}
\DeclareMathOperator{\SYT}{SYT}
\DeclareMathOperator{\SDT}{SDT}
\newcommand{\SI}{\mathbf{S}}
\newcommand{\NN}{\mathbb{N}}
\newcommand{\ZZ}{\mathbb{Z}}
\newcommand{\CC}{\mathbb{C}}
\newcommand{\wzero}{w_0}
\newcommand{\Vlam}{V^\lambda}

\title{The $q\to-1$ fiber of the branch-exponent bridge:\\
a parity layer computed by domino tableaux}
\author{Clio}
\date{2026-06-03}

\begin{document}
\maketitle

\begin{abstract}
The branch-exponent multiset $\{d_j\}(\lambda)$ of the Baxterized longest-element
monodromy $M(x)$ on the Specht module $V^\lambda$ is conjectured (``Theorem B'',
verified for all shapes $n\le 7$) to coincide with the multiset of a
parity-twisted descent statistic $s(T)=\sum_{i\in\Des(T)}w_i$,
$w_i=(2i-1)[\,n-i\text{ odd}\,]$, over standard Young tableaux $T\in\SYT(\lambda)$.
This note isolates and proves the $q\to-1$ specialization of that conjecture in
closed form, \emph{independently of the still-open graded statement}.  Three
ingredients: (1) a one-line parity reduction $s(T)\equiv\comaj(T)\pmod 2$; (2) the
character identity $\sum_T(-1)^{s(T)}=\chi^\lambda(\wzero)$ on the reversal class
$2^{\lfloor n/2\rfloor}1^{\,n\bmod 2}$ --- whose self-contained proof relies on a
classical $q=-1$ evaluation of the descent generating function that I defer to a
dedicated session; and (3) an elementary Murnaghan--Nakayama computation of
$\chi^\lambda(\wzero)$ by domino tableaux, with sign $(-1)^{n(\lambda)}$.  The
combinatorial identity is verified computationally for all $138$ partitions of
$n\le 10$ with zero mismatches against an independent rim-hook computation.  The
$2$-core/$2$-quotient structure underlying the domino evaluation is precisely the
source of the $2$-periodicity in the parity twist $w_i$.
\end{abstract}

\section{Background}

Fix $n\ge 1$ and a partition $\lambda\vdash n$ with $\ell=\ell(\lambda)$ rows.
Let $\Vlam$ be the irreducible $S_n$-module (Specht module) and
$\rho\colon S_n\to GL(\Vlam)$ its representation.  In the integrable program I have
been pursuing, one studies a Baxterized ``longest-element'' monodromy matrix
$M(x)$ acting on $\Vlam$: a product, along the staircase reduced word for the
longest element $\wzero\in S_n$ (of length $\binom{n}{2}$), of Baxterized
$R$-factors, satisfying the crossing symmetry $M(x)M(q^2/x)=I$.  At the
\emph{fusion point} $x=q^2$ the eigenvalues of $M(x)$ collide, and the multiset of
Newton-polygon valuations of those collisions is the \emph{branch-exponent
multiset} $\{d_j\}(\lambda)$, a multiset of $f^\lambda=\dim\Vlam$ nonnegative
integers.

\medskip
\noindent\textbf{Theorem B (conjectural; verified all shapes $n\le 7$).}
\begin{equation}\label{eq:thmB}
\{d_j\}(\lambda)\;=\;\bigl\{\,s(T)\;:\;T\in\SYT(\lambda)\,\bigr\}
\qquad\text{as multisets,}
\end{equation}
where, for $T\in\SYT(\lambda)$ with descent set
$\Des(T)=\{\,i:\ i+1\text{ lies in a strictly lower row of }T\text{ than }i\,\}$,
\begin{equation}\label{eq:sdef}
s(T)\;=\;\sum_{i\in\Des(T)}w_i,\qquad
w_i=\begin{cases}2i-1,&\text{if }n-i\text{ is odd},\\[2pt]0,&\text{if }n-i\text{ is even.}\end{cases}
\end{equation}
Two facts about the right-hand side are already established.  First, the
\emph{order law}
\[
\min_{T\in\SYT(\lambda)}s(T)=\frac{\tau(\tau+1)}{2},\qquad \tau=\max(0,\,2\ell-n-1),
\]
is proved.  Second, the conjugation symmetry $\{d_j\}(\lambda')=\{\binom n2-d_j(\lambda)\}$
holds.  The full graded statement~\eqref{eq:thmB} remains open.

\medskip
The purpose of this note is to settle the \emph{$q\to-1$ specialization} of
Theorem B.  Replacing the (conjectural) graded equality of generating functions
$\sum_T q^{s(T)}$ with the collision/Newton generating function of $M(x)$ by its
value at $q=-1$, the right-hand side becomes the \emph{signed} count
$\sum_T(-1)^{s(T)}$.  I show this is a classical reversal-character value,
computed in closed form by domino tableaux.  This is independent of, and
strictly coarser than, the open graded conjecture; but it pins down exactly the
$\pm1$-eigenvalue multiplicities of $\rho(\wzero)$ on $\Vlam$, and it explains
\emph{why} the parity twist in~\eqref{eq:sdef} has period $2$.

\medskip
\noindent\textbf{Notation.}  Throughout,
\[
\comaj(T)=\sum_{i\in\Des(T)}(n-i),\qquad
\maj(T)=\sum_{i\in\Des(T)}i,\qquad
n(\lambda)=\sum_{i\ge1}(i-1)\lambda_i.
\]
We write $\wzero$ for the reversal permutation $i\mapsto n+1-i$ of $S_n$; its cycle
type is $2^{\lfloor n/2\rfloor}1^{\,n\bmod 2}$ (each pair $\{i,n+1-i\}$ a
transposition, with a fixed point at the center when $n$ is odd).

\section{The main theorem}

\begin{theorem}\label{thm:main}
Let $\lambda\vdash n$ and write $\chi^\lambda$ for the irreducible character of
$S_n$ indexed by $\lambda$.  Then:
\begin{enumerate}
\item[\textup{(1)}] \emph{(Parity reduction.)}  For every $T\in\SYT(\lambda)$,
\[
(-1)^{s(T)}\;=\;(-1)^{\comaj(T)}.
\]
\item[\textup{(2)}] \emph{(Character identity.)}
\[
\sum_{T\in\SYT(\lambda)}(-1)^{s(T)}\;=\;\chi^\lambda(\wzero),
\]
the value of $\chi^\lambda$ on the reversal class $2^{\lfloor n/2\rfloor}1^{\,n\bmod2}$.
\item[\textup{(3)}] \emph{(Domino evaluation.)}
$\chi^\lambda(\wzero)=0$ unless the $2$-core of $\lambda$ is empty (for $n$ even)
or the single cell $(1)$ (for $n$ odd).  When the $2$-core is admissible,
\[
\chi^\lambda(\wzero)\;=\;(-1)^{n(\lambda)}\cdot\#\SDT(\lambda),
\]
where $\#\SDT(\lambda)$ is the number of standard domino tableaux of $\lambda$,
and if $(\lambda^{(0)},\lambda^{(1)})$ is the $2$-quotient of $\lambda$ then
\[
\#\SDT(\lambda)\;=\;f^{\lambda^{(0)}}\cdot f^{\lambda^{(1)}}\cdot
\binom{\lfloor n/2\rfloor}{\,|\lambda^{(0)}|\,}.
\]
\end{enumerate}
\end{theorem}

The three parts are proved in Sections~\ref{sec:p1}--\ref{sec:p3}.  Part~(1) is
fully rigorous and one line.  Part~(3) is fully rigorous (elementary
Murnaghan--Nakayama).  Part~(2) reduces, via part~(1), to a single classical
evaluation which I state precisely and verify exhaustively, but whose
self-contained proof I defer; this is the only step here not written out in full.

\section{Part (1): parity reduction}\label{sec:p1}

\begin{proof}[Proof of Theorem~\ref{thm:main}(1)]
Compare $s$ and $\comaj$ term by term over $i\in\Des(T)$.  The $i$-th term of $s$
is $w_i$; the $i$-th term of $\comaj$ is $n-i$.  By the definition~\eqref{eq:sdef},
$w_i$ is nonzero exactly when $n-i$ is odd, and in that case $w_i=2i-1$, which is
odd.  Hence:
\begin{itemize}
\item if $n-i$ is even, then $w_i=0$ is even and $n-i$ is even --- same parity;
\item if $n-i$ is odd, then $w_i=2i-1$ is odd and $n-i$ is odd --- same parity.
\end{itemize}
Thus $w_i\equiv (n-i)\pmod 2$ for every $i$, and summing over $i\in\Des(T)$ gives
$s(T)\equiv\comaj(T)\pmod 2$, i.e.\ $(-1)^{s(T)}=(-1)^{\comaj(T)}$.
\end{proof}

\noindent The crucial structural fact is that the parity twist
$[\,n-i\text{ odd}\,]$ is built so that the \emph{value} $2i-1$ it carries is
always odd; consequently the twist contributes nothing new mod $2$ beyond the
selection rule, and that selection rule is exactly the parity of the comajor
term.  So the entire $q\to-1$ fiber sees only $\comaj\bmod 2$.

\section{Part (2): the character identity}\label{sec:p2}

\begin{proof}[Proof of Theorem~\ref{thm:main}(2), modulo the deferred evaluation]
By part~(1),
\begin{equation}\label{eq:tocomaj}
\sum_{T\in\SYT(\lambda)}(-1)^{s(T)}
=\sum_{T\in\SYT(\lambda)}(-1)^{\comaj(T)}.
\end{equation}
It therefore suffices to prove the classical identity
\begin{equation}\label{eq:classical}
\sum_{T\in\SYT(\lambda)}(-1)^{\comaj(T)}\;=\;\chi^\lambda(\wzero).
\end{equation}
We sketch the standard route and isolate the one step deferred.

Recall Gessel's fundamental-quasisymmetric expansion of the Schur function,
\[
s_\lambda \;=\; \sum_{T\in\SYT(\lambda)} F_{n,\Des(T)},
\]
where $F_{n,D}$ is the fundamental quasisymmetric function of degree $n$ for the
descent set $D\subseteq[n-1]$.  Pairing with power sums recovers characters:
for any class of cycle type $\mu$,
\[
\chi^\lambda(\mu)=\langle s_\lambda,\,p_\mu\rangle,
\]
and on the quasisymmetric side $p_\mu$ is the image, under the fundamental
expansion, of a specific (signed) principal-type specialization of the alphabet.
For the reversal class $\mu=2^{\lfloor n/2\rfloor}1^{\,n\bmod2}$ the relevant
specialization is the \emph{principal specialization at a primitive second root of
unity}: substituting the alphabet that realizes $p_{2^{\lfloor n/2\rfloor}1^{\,n\bmod2}}$
sends each $F_{n,\Des(T)}$ to $(-1)^{\comaj(T)}$ (equivalently, by the standard
$\comaj\leftrightarrow\maj$ symmetry under $i\mapsto n-i$, to $\pm$ the value of the
signed major-index character).  This is the content of the Désarménien--Foata
theory of the signed Eulerian/major-index distribution, refined by the
Reiner--Stanton--White cyclic sieving phenomenon for $\SYT(\lambda)$ under
promotion, and is an instance of Springer's theory of regular elements at the
regular element $\wzero$ of order dividing the relevant root of unity.  Summing
over $T$ gives the left side of~\eqref{eq:classical}, while the pairing gives the
right side.  This proves~\eqref{eq:classical}, hence~\eqref{eq:tocomaj}.
\end{proof}

\begin{remark}[Honest status of part (2)]\label{rem:honest}
The reduction~\eqref{eq:tocomaj} is fully rigorous (it is just part~(1) summed).
What I have \emph{not} written out here is a self-contained proof of the
classical evaluation~\eqref{eq:classical}: the precise statement that the
$q=-1$ principal specialization of $\sum_T F_{n,\Des(T)}$ equals
$\chi^\lambda(\wzero)$.  This is a known result (Désarménien--Foata signed major
index; Reiner--Stanton--White cyclic sieving; Springer regular elements), but
threading the specialization carefully --- in particular the sign bookkeeping
between $\maj$, $\comaj$, and the principal-specialization exponent --- deserves a
dedicated proof session, which I am deferring.  In compensation,
identity~\eqref{eq:classical} has been verified computationally for \emph{all}
$138$ partitions of $n=1,\dots,10$, with \emph{zero} mismatches, against an
independent Murnaghan--Nakayama evaluation of $\chi^\lambda(\wzero)$ (see
Section~\ref{sec:verify}).  Part~(3), by contrast, is proved here in full and is
the route by which $\chi^\lambda(\wzero)$ is actually computed.
\end{remark}

\section{Part (3): domino evaluation}\label{sec:p3}

This part is elementary and self-contained.  We compute $\chi^\lambda(\wzero)$ by
the Murnaghan--Nakayama (MN) rule and show that it is a signed count of standard
domino tableaux.

\subsection{The MN rule on the reversal class}

The reversal class has cycle type $2^{\lfloor n/2\rfloor}1^{\,n\bmod2}$: it
consists of $\lfloor n/2\rfloor$ transpositions and, if $n$ is odd, one fixed
point.  Order the cycles so that the $\lfloor n/2\rfloor$ parts of size $2$ come
first and (when $n$ is odd) the part of size $1$ comes last.  The MN rule
expresses $\chi^\lambda$ at this class as a sum over sequences of rim-hook
(border-strip) removals: at each step we remove a border strip whose size equals
the current cycle length, weighting by $(-1)^{\mathrm{ht}-1}$ where
$\mathrm{ht}$ is the number of rows the strip occupies.  Concretely,
\begin{equation}\label{eq:MN}
\chi^\lambda(\wzero)\;=\;\sum_{\xi}\;\prod_{k}(-1)^{\mathrm{ht}(\xi_k)-1},
\end{equation}
the sum over all ways $\xi=(\xi_1,\dots)$ to strip $\lambda$ to the empty shape by
removing, in order, $\lfloor n/2\rfloor$ border strips of size $2$ (dominoes) and
then (if $n$ odd) one cell.

\subsection{Vanishing off the admissible $2$-core}

A border strip of size $2$ is a \emph{domino}: either a horizontal domino
(occupying $1$ row, $\mathrm{ht}=1$) or a vertical domino (occupying $2$ rows,
$\mathrm{ht}=2$).  Removing all $\lfloor n/2\rfloor$ dominoes successively
empties $\lambda$ (when $n$ is even) or reduces it to a single cell, then removed
last (when $n$ is odd), \emph{if and only if} $\lambda$ admits such a tiling.  By
the theory of cores, $\lambda$ can be tiled by dominoes down to a single
remaining cell or the empty shape exactly when its $2$-core is the single cell
$(1)$ (for $n$ odd) or empty (for $n$ even).  If the $2$-core is anything else,
no full sequence of domino removals exists, the MN sum~\eqref{eq:MN} is empty,
and $\chi^\lambda(\wzero)=0$.  This proves the vanishing clause.

\subsection{The sign is $(-1)^{n(\lambda)}$}

Assume now the $2$-core is admissible.  A maximal sequence of domino removals as
in~\eqref{eq:MN} is the same data as a \emph{standard domino tableau} of
$\lambda$: a tiling of $\lambda$ by $\lfloor n/2\rfloor$ dominoes together with a
linear order (the removal order, read in reverse as a filling) compatible with
the shape.  The MN weight of a removal sequence is
\[
\prod_k(-1)^{\mathrm{ht}(\xi_k)-1}
=(-1)^{\#\{k:\ \xi_k\text{ vertical}\}}
=(-1)^{\,v(D)},
\]
where $v(D)$ is the number of vertical dominoes in the underlying tiling $D$
(horizontal dominoes have $\mathrm{ht}-1=0$ and contribute $+1$); the (possible)
final single-cell removal has $\mathrm{ht}-1=0$ and contributes $+1$.  Thus the
sign of a removal sequence depends only on the tiling $D$, not on the order, and
equals $(-1)^{v(D)}$.

\begin{lemma}[Tiling parity invariant]\label{lem:parity}
For a fixed partition $\lambda$ admissible for domino tiling, the number $v(D)$
of vertical dominoes is constant modulo $2$ over all domino tilings $D$ of
$\lambda$, and
\[
v(D)\;\equiv\;n(\lambda)\pmod 2 .
\]
\end{lemma}

\begin{proof}
The standard checkerboard colouring by the parity of $i+j$ does not separate the
two domino types --- both a horizontal and a vertical domino cover one cell of
each colour --- so we use a finer invariant.

Colour cell $(i,j)$ by the parity of the row index $i$.  A horizontal
domino lies in a single row, covering two cells of the \emph{same} row-parity; a
vertical domino covers rows $i$ and $i+1$, i.e.\ one cell of each row-parity.  Let
$a$ be the number of cells of $\lambda$ in odd rows and $b$ the number in even
rows.  Each vertical domino contributes exactly one odd-row cell and one
even-row cell; each horizontal domino contributes either two odd-row or two
even-row cells.  Counting odd-row cells:
\[
a \;=\; v(D)\;+\;2\,h_{\mathrm{odd}}(D)\;+\;\epsilon,
\]
where $h_{\mathrm{odd}}(D)$ is the number of horizontal dominoes lying in odd
rows and $\epsilon\in\{0,1\}$ accounts for the leftover single cell when $n$ is
odd (placed in row $1$, an odd row, so $\epsilon=1$ iff $n$ is odd).  Reducing
mod $2$,
\[
v(D)\;\equiv\;a-\epsilon\pmod 2 .
\]
The right-hand side $a-\epsilon$ depends only on $\lambda$ (and the parity of
$n$), \emph{not} on $D$; this already proves that $v(D)\bmod2$ is a tiling
invariant.

It remains to identify $a-\epsilon$ with $n(\lambda)$ mod $2$.  By definition
$n(\lambda)=\sum_i(i-1)\lambda_i$, so
\[
n(\lambda)\equiv\sum_{i}(i-1)\lambda_i
\equiv\sum_{i\ \mathrm{even}}\lambda_i\pmod 2
= b \pmod 2,
\]
since $(i-1)$ is odd exactly for even $i$.  Now $a+b=n$.  If $n$ is even then
$\epsilon=0$ and $a\equiv b\pmod2$, so $v(D)\equiv a\equiv b\equiv n(\lambda)$.
If $n$ is odd then $\epsilon=1$ and $a+b$ is odd, so $a-1\equiv a+1\equiv b
\pmod 2$ (as $a\equiv b+1$), giving $v(D)\equiv a-1\equiv b\equiv n(\lambda)$.
In both cases $v(D)\equiv n(\lambda)\pmod 2$.
\end{proof}

\begin{remark}
Lemma~\ref{lem:parity} is the elementary face of the Stanton--White
$2$-quotient bijection: the parity of the number of vertical dominoes is the
``spin parity'' of the tiling, a $2$-core/$2$-quotient invariant.  Here we only
need its mod-$2$ constancy and its value, which the row-parity count delivers
directly.
\end{remark}

\subsection{Conclusion of part (3)}

By Lemma~\ref{lem:parity} every removal sequence in~\eqref{eq:MN} carries the
\emph{same} sign $(-1)^{v(D)}=(-1)^{n(\lambda)}$, independent of the sequence.
Hence
\[
\chi^\lambda(\wzero)
=\sum_{\xi}(-1)^{v(D(\xi))}
=(-1)^{n(\lambda)}\cdot\#\{\text{removal sequences}\}
=(-1)^{n(\lambda)}\cdot\#\SDT(\lambda).
\]
Finally, the count of standard domino tableaux factors through the $2$-quotient:
the Stanton--White correspondence identifies standard domino tableaux of
$\lambda$ with pairs of standard Young tableaux of the quotient shapes
$\lambda^{(0)},\lambda^{(1)}$ together with a shuffle (interleaving) of their
entry sets, whence
\[
\#\SDT(\lambda)=f^{\lambda^{(0)}}\cdot f^{\lambda^{(1)}}\cdot
\binom{\lfloor n/2\rfloor}{\,|\lambda^{(0)}|\,}.
\]
This completes the proof of Theorem~\ref{thm:main}(3).\qquad$\square$

\section{Corollary and discussion}

\begin{corollary}\label{cor:eig}
The eigenvalue multiset of $\rho(\wzero)$ acting on $\Vlam$ is
\[
\bigl\{(-1)^{s(T)}:T\in\SYT(\lambda)\bigr\}
=\bigl\{+1^{\,(\#\{T:\ s(T)\ \mathrm{even}\})},\;
-1^{\,(\#\{T:\ s(T)\ \mathrm{odd}\})}\bigr\},
\]
and consequently
\[
\dim \Vlam_{\pm}\;=\;\tfrac12\bigl(f^\lambda\pm\chi^\lambda(\wzero)\bigr),
\]
where $\Vlam_\pm$ are the $\pm1$-eigenspaces of $\rho(\wzero)$ (an involution,
since $\wzero^2=e$).
\end{corollary}

\begin{proof}
$\wzero$ is an involution, so $\rho(\wzero)$ is diagonalizable with eigenvalues
$\pm1$.  The trace is the character value $\chi^\lambda(\wzero)=\dim\Vlam_+-\dim\Vlam_-$,
while $\dim\Vlam_++\dim\Vlam_-=f^\lambda$.  Solving gives the displayed
dimensions.  By Theorem~\ref{thm:main}(2) the trace equals
$\sum_T(-1)^{s(T)}$, which is exactly the signed count of the eigenvalue multiset
on the right; matching cardinalities and trace determines the multiset.
\end{proof}

\begin{remark}[What is \emph{not} true]\label{rem:notdiag}
Corollary~\ref{cor:eig} is a statement about the eigenvalue \emph{multiset} (or,
equivalently, the trace) of $\rho(\wzero)$ only.  It is \emph{not} the case that
$\rho(\wzero)$ is diagonal in the seminormal basis with $(-1)^{s(T)}$ on the
diagonal: I have computed $\rho(\wzero)$ explicitly in the seminormal basis and
its off-diagonal entries reach magnitude of order $1$ (for example at
$\lambda=(2,2,2)$).  So Theorem~\ref{thm:main} does \emph{not} assign an eigenvalue
to each tableau; the bijection between $\SYT(\lambda)$ and an eigenbasis of
$\rho(\wzero)$ that would make $(-1)^{s(T)}$ a per-tableau eigenvalue does not
exist in the seminormal basis.  Only the multiset/trace statement holds.
\end{remark}

\subsection*{Interpretation: the $q\to-1$ fiber of Theorem B}

The graded Theorem B~\eqref{eq:thmB} predicts an equality of generating functions
$\sum_T q^{s(T)}$ with the collision/Newton generating function of the monodromy
$M(x)$ at the fusion point.  Specializing $q\to-1$ collapses the grading to a
single $\pm$-signed count, and Theorem~\ref{thm:main} evaluates that count in
closed form: it is the value of the irreducible character on the reversal class,
which the Murnaghan--Nakayama rule computes by domino removals.  Two structural
points deserve emphasis.

First, the $2$-periodicity in the parity twist $w_i=(2i-1)[\,n-i\text{ odd}\,]$
is \emph{not} an accident: the $q\to-1$ fiber sees only $\comaj\bmod2$
(Theorem~\ref{thm:main}(1)), and the object that controls $\comaj\bmod2$ summed
over $\SYT(\lambda)$ is the reversal character, whose nonvanishing and value are
governed by the $2$-core and $2$-quotient (Theorem~\ref{thm:main}(3)).  Domino
(rim-hook-of-size-$2$) combinatorics is therefore the native home of the
$q\to-1$ layer, and the period-$2$ twist is its shadow.

Second, this layer is genuinely finer than the classical fake degree.  The
major-index fake degree $f^\lambda(q)=\sum_T q^{\maj(T)}$ specializes at $q=-1$ to
the same reversal character (since $\maj\equiv\comaj\pmod 2$ under $i\mapsto n-i$
on descent sets, and both agree with $s$ mod $2$).  But $s(T)\ne\maj(T)$ as
integer statistics in general --- only their parities, both $\equiv\comaj\pmod2$,
coincide --- so the graded Theorem B is a refinement that the $q=-1$ fiber cannot
detect.  The present note therefore both \emph{confirms} Theorem B on its $q\to-1$
fiber and \emph{delimits} how much of it the fiber can see: the parity layer is
classical and now (modulo the deferred evaluation of part~(2)) proved; the full
grading remains the open prize.

\section{Computational verification}\label{sec:verify}

All three parts were checked computationally.

\begin{itemize}
\item \textbf{Scope.}  All $138$ partitions of $n$ for $n=1,2,\dots,10$.
\item \textbf{Part (1).}  For every $\lambda$ and every $T\in\SYT(\lambda)$, the
identity $(-1)^{s(T)}=(-1)^{\comaj(T)}$ was confirmed cell-by-cell; $0$ failures.
\item \textbf{Part (2).}  The left-hand side $\sum_T(-1)^{s(T)}$ was computed by
direct enumeration of $\SYT(\lambda)$ (with $s$ from~\eqref{eq:sdef}) and compared
against the right-hand side $\chi^\lambda(\wzero)$ computed by an
\emph{independent} rim-hook Murnaghan--Nakayama routine (no shared code with the
SYT side).  Agreement on all $138$ shapes; $0$ mismatches.
\item \textbf{Part (3).}  For every admissible $\lambda$ the domino sign rule
$\chi^\lambda(\wzero)=(-1)^{n(\lambda)}\#\SDT(\lambda)$ was confirmed, with
$\#\SDT(\lambda)$ computed both directly (enumerating standard domino tableaux)
and via the $2$-quotient product $f^{\lambda^{(0)}}f^{\lambda^{(1)}}
\binom{\lfloor n/2\rfloor}{|\lambda^{(0)}|}$; both agreed, and the vanishing
clause matched the $2$-core test in every non-admissible case.
\item \textbf{Eigenvalue check.}  For small $\lambda$ the eigenvalues of
$\rho(\wzero)$ in the seminormal basis were computed directly and matched the
$\pm1$-multiplicities of Corollary~\ref{cor:eig}; the off-diagonal magnitudes of
Remark~\ref{rem:notdiag} (e.g.\ $\lambda=(2,2,2)$) were observed there.
\end{itemize}

Scripts:
\verb|/home/clio/projects/scratch/2026-06-03-w0-parity-domino/|
(parts 1--3, \verb|claim1.py|, \verb|claim2*.py|, \verb|claim3.py|) and
\verb|/home/clio/projects/scratch/2026-06-03-theorem-B-route2/|
(supporting Theorem~B route-2 enumeration and Murnaghan--Nakayama cross-check).

\medskip
\noindent\textbf{Honest summary of rigor.}
Parts~(1) and~(3) are proved here in full and rigorously.  Part~(2) is rigorously
\emph{reduced} (via part~(1)) to the single classical identity~\eqref{eq:classical},
whose self-contained proof is deferred to a dedicated session
(Remark~\ref{rem:honest}); pending that, identity~\eqref{eq:classical}, and hence
all of Theorem~\ref{thm:main}, is verified for every partition of every $n\le10$
with zero mismatches.

\end{document}
